\documentclass[11pt,a4paper]{article}

% Packages
\usepackage[utf8]{inputenc}
\usepackage[T1]{fontenc}
\usepackage{amsmath,amssymb}
\usepackage{graphicx}
\usepackage{booktabs}
\usepackage{multirow}
\usepackage{array}
\usepackage{float}
\usepackage{hyperref}
\usepackage{natbib}
\usepackage{geometry}
\usepackage{xcolor}
\usepackage{algorithm}
\usepackage{algpseudocode}

% Page geometry
\geometry{margin=2.5cm}

% Line spacing
\usepackage{setspace}
\onehalfspacing

% Title and author
\title{\textbf{PlantHGNN: Integrating Heterogeneous Biological Networks via Multi-View Graph Neural Networks for Plant Genomic Prediction}}
\author{Guohao Lyu$^{1,2}$, Yingchun Xia$^{1}$, Huichao Liu$^{1}$, Xiaolei Zhu$^{1}$, Shuai Yang$^{1}$, Qingyong Wang$^{1,*}$, Lichuan Gu$^{1,2,*}$\\
\small{$^1$School of Artificial Intelligence, Anhui Agricultural University}\\
\small{$^2$Key Laboratory of Smart Agriculture Technology and Equipment of Anhui Province}\\
\small{*Corresponding authors: glc@ahau.edu.cn}}
\date{}

\begin{document}

\maketitle

% Abstract
\begin{abstract}
Genomic prediction (GP) is essential for accelerating plant breeding programs. Although Graph Neural Networks (GNNs) have shown potential in GP tasks, effectively integrating multiple heterogeneous biological networks remains challenging. In this study, we propose PlantHGNN, a multi-view graph neural network framework that adaptively fuses Protein-Protein Interaction (PPI) and KEGG pathway networks for plant genome prediction. Through comprehensive experiments on the rice GSTP007 dataset, our multi-view approach achieves a Pearson correlation coefficient (PCC) of 0.8802 on the Grain\_Length trait, outperforming single-view NetGP (0.8760) and DNNGP (0.8751). Systematic ablation experiments demonstrate that PPI+KEGG integration (PCC=0.8608) outperforms individual networks (PPI: 0.8541; KEGG: 0.8498), with learned attention weights confirming the complementary contributions of each view. Five-fold cross-validation across three random seeds (15 total runs) confirms the stability of our results (MultiView: 0.880$\pm$0.015 versus NetGP: 0.876$\pm$0.016). Multi-trait validation demonstrates competitive performance across all six agronomic traits. These findings support the effectiveness of multi-view heterogeneous network integration for plant genome prediction and provide a framework for leveraging multiple biological networks in genomic selection.

Keywords: Genomic prediction, Graph neural networks, Multi-view learning, Plant breeding, Heterogeneous networks
\end{abstract}

% Introduction
\section{Introduction}
\label{sec:introduction}

Genomic prediction (GP) enables the prediction of complex traits from genome-wide markers and has become an essential tool in modern plant breeding \citep{meuwissen2001prediction}. By estimating genomic estimated breeding values (GEBVs) from high-density single nucleotide polymorphism (SNP) markers, GP accelerates the selection cycle and improves genetic gain compared to traditional phenotypic selection \citep{crossa2017genomic}. Traditional statistical methods such as GBLUP (Genomic Best Linear Unbiased Prediction) and Bayesian approaches have been widely adopted in breeding programs due to their computational efficiency and interpretability \citep{vanraden2008efficient,hayes2019100}. However, these approaches often encounter limitations when modeling the complex non-linear relationships between genetic markers and phenotypic traits, particularly for polygenic traits influenced by numerous small-effect variants \citep{meuwissen2001prediction}.

The emergence of deep learning has created new opportunities for improving the accuracy of GP. Deep Neural Networks (DNNs) \citep{zhang2018deep}, Convolutional Neural Networks (CNNs) \citep{ma2021cropformer}, Recurrent Neural Networks (RNNs) \citep{wang2023deep}, and Graph Neural Networks (GNNs) \citep{wang2021netgp} have all been applied to genomic prediction with varying degrees of success. These methods can automatically learn hierarchical representations from raw genetic data, potentially capturing epistatic interactions that traditional methods might miss \citep{montesinos2021application}. Among these approaches, GNNs are particularly promising because of their natural ability to incorporate biological network information, including protein-protein interaction (PPI) networks, Gene Ontology (GO) annotations, and pathway data \citep{qiu2023netgp,zhao2025netgp}.

Despite the growing application of GNNs in genomic prediction, several critical challenges remain. First, most existing GNN-based approaches rely on a single type of biological network (typically PPI), leaving abundant functional information from pathway and ontology databases underutilized \citep{wang2021netgp,qiu2023netgp}. Second, the integration of multiple heterogeneous networks in plant GP remains underexplored. While multi-view learning has shown promise in other domains, systematic investigations of how different network combinations affect prediction accuracy in plant GP are scarce \citep{hu2025hgnn}. Third, rigorous statistical validation is often lacking, with many studies reporting single-point estimates without confidence intervals, significance testing, or ablation experiments to isolate the contributions of specific methodological choices.

In this study, we address these gaps through a comprehensive investigation of PlantHGNN, a multi-view heterogeneous graph neural network framework for plant multi-trait genomic prediction. Our contributions are fourfold:

\begin{enumerate}
    \item \textbf{Multi-View Heterogeneous Network Integration}: We present PlantHGNN, a framework that adaptively integrates multiple biological networks (PPI, GO, KEGG) through attention-based fusion. To our knowledge, this represents one of the first systematic evaluations of heterogeneous network integration for plant genomic prediction.

    \item \textbf{Rigorous Ablation and Component Analysis}: We conduct comprehensive ablation experiments isolating the contributions of each network view and architectural component. Our results demonstrate that PPI+KEGG integration achieves the best performance (PCC=0.8608), outperforming individual networks (PPI: 0.8541; KEGG: 0.8498) and revealing insights into network complementarity and redundancy.

    \item \textbf{Statistically Robust Validation}: We validate our findings through 5-fold cross-validation across multiple random seeds (15 runs per comparison), paired with comprehensive statistical testing including paired t-tests and Wilcoxon signed-rank tests with Bonferroni correction. This provides robust, reproducible estimates of performance differences.

    \item \textbf{Multi-Trait Generalization and Biological Interpretability}: We extend validation to all six agronomic traits in the GSTP007 dataset and analyze the learned attention weights to understand which network views contribute most to different trait types, providing biologically grounded insights into the prediction mechanisms.
\end{enumerate}

Through these contributions, we advance the understanding of how heterogeneous biological network integration can improve GNN-based genomic prediction. Our findings offer actionable insights for practitioners seeking to deploy these methods in real-world breeding programs, while our open-source implementation enables reproducibility and extension by the research community.

% Methods

\section{Materials and Methods}
\label{sec:methods}

\subsection{Dataset}
\label{subsec:dataset}

We conducted experiments on the GSTP007 rice dataset. The dataset comprises 1,269 rice accessions with 5,291 high-quality single nucleotide polymorphisms. Six agronomic traits were measured: Grain\_Length, Grain\_Width, Grain\_Weight, Panicle\_Length, Plant\_Height, and Yield\_per\_plant. For network-based modeling, we employed a Protein-Protein Interaction (PPI) network containing 831 nodes and a KEGG pathway network with 4,053 nodes.

The dataset was divided into training (80\%, n=1,015) and validation (20\%, n=254) sets using a random seed of 42 to ensure reproducibility.

\subsection{Model Architecture}
\label{subsec:model}

Figure~\ref{fig:architecture} illustrates the overall architecture of PlantHGNN. Our model integrates multiple biological networks through a multi-view learning framework.

\begin{figure*}[t]
\centering
\includegraphics[width=0.95\textwidth]{../figures/figure1_architecture.png}
\caption{PlantHGNN architecture. The model consists of four main components: (A) Multi-view GCN encoders process PPI and KEGG networks in parallel as separate homogeneous graphs; (B) Attention-based fusion dynamically weights each network's contribution; (C) Functional and structural encoding incorporates biological knowledge; (D) AttnRes Transformer encoder captures deep feature interactions for final GEBV prediction.}
\label{fig:architecture}
\end{figure*}

PlantHGNN is a multi-view graph neural network that integrates four main components. Multi-view GCN encoders process the PPI and KEGG networks through separate Graph Convolutional Networks. An attention-based fusion module combines multi-view representations using learnable attention weights. The AttnRes Transformer employs attention residual blocks for deep feature extraction. Finally, a regression head produces predictions of genomic estimated breeding values (GEBVs).

\subsection{Hyperparameter Optimization}
\label{subsec:hpo}

We systematically optimized the following hyperparameters:

\begin{table}[H]
\centering
\caption{Hyperparameter search space}
\label{tab:hpo_space}
\begin{tabular}{lll}
\toprule
\textbf{Parameter} & \textbf{Search Range} & \textbf{Description} \\
\midrule
d\_model & \{64, 128, 256\} & Model embedding dimension \\
batch\_size & \{32, 64, 128\} & Training batch size \\
learning\_rate & \{1e-4, 1.4e-4, 5e-4\} & Initial learning rate \\
dropout & \{0.20, 0.25, 0.30\} & Dropout rate \\
\bottomrule
\end{tabular}
\end{table}

The complete search space contains $3 \times 3 \times 3 \times 3 = 81$ possible configurations. Due to computational constraints, we evaluated 8 configurations (9.9\% coverage), prioritizing combinations expected to perform well based on preliminary experiments. We acknowledge that this limited coverage may miss the global optimum; automated methods such as Bayesian optimization would explore the space more efficiently (see Limitations).

\section{Results}
\label{sec:results}

\subsection{Model Configuration via Hyperparameter Optimization}
\label{subsec:hpo_results}

To ensure fair evaluation of the multi-view architecture, we first identified a strong base configuration through limited hyperparameter optimization on the Grain\_Length trait. Table~\ref{tab:hpo_results} summarizes the configurations tested.

\begin{table}[H]
\centering
\caption{Hyperparameter search results on Grain\_Length trait}
\label{tab:hpo_results}
\begin{tabular}{cccccc}
\toprule
\textbf{Config} & \textbf{d\_model} & \textbf{Batch} & \textbf{LR} & \textbf{Dropout} & \textbf{PCC} \\
\midrule
Baseline & 64 & 32 & 5e-4 & 0.20 & 0.8610 \\
HPO-1 & 128 & 64 & 1.4e-4 & 0.25 & 0.8686 \\
HPO-2 & 256 & 64 & 1.4e-4 & 0.25 & 0.8704 \\
HPO-3 & 128 & 128 & 1.4e-4 & 0.25 & 0.7575 \\
\bottomrule
\end{tabular}
\end{table}

The HPO-2 configuration (d\_model=256, batch\_size=64, lr=1.4e-4, dropout=0.25) achieved the best validation PCC of 0.8704 and was used as the base configuration for all subsequent multi-view and ablation experiments (Figure~\ref{fig:hpo_comparison}).

\begin{figure}[H]
\centering
\includegraphics[width=0.8\textwidth]{../figures/figure1_hpo_comparison.pdf}
\caption{Hyperparameter optimization results. HPO-2 (d\_model=256) achieves the best PCC of 0.8704, with 1.09\% improvement over baseline.}
\label{fig:hpo_comparison}
\end{figure}

\subsection{Cross-Validation Results}
\label{subsec:cv_results}

To validate the stability of our results, we conducted 5-fold cross-validation using the optimal HPO configuration. Table~\ref{tab:cv_results} presents the detailed fold-by-fold results comparing GBLUP, Ridge Regression, DNNGP, NetGP, and our MultiView approach.

\input{tables/table_5fold_cv_grain_length}

The 5-fold CV results show consistent performance across folds. The MultiView\_PPI\_GO configuration achieves a mean PCC of 0.880$\pm$0.015 across 15 runs, which exceeds the performance of NetGP (0.876$\pm$0.016) and DNNGP (0.875$\pm$0.015). Table~\ref{tab:stats_test} presents comprehensive statistical significance tests including paired t-tests and Wilcoxon signed-rank tests with 95\% confidence intervals.

\begin{table}[H]
\centering
\caption{Statistical significance tests (5-fold CV $\times$ 3 seeds, n=15)}
\label{tab:stats_test}
\begin{tabular}{lccc}
\toprule
\textbf{Comparison} \& \textbf{Test} \& \textbf{Statistic} \& \textbf{$p$-value} \\
\midrule
MultiView vs GBLUP \& Paired t-test \& $t = 8.94$ \& $<$0.001$^{***}$ \\
MultiView vs GBLUP \& Wilcoxon signed-rank \& $W = 120$ \& $<$0.001$^{***}$ \\
MultiView vs DNNGP \& Paired t-test \& $t = 2.41$ \& 0.031$^{*}$ \\
MultiView vs NetGP \& Paired t-test \& $t = 1.85$ \& 0.084 (ns) \\
MultiView vs NetGP \& Wilcoxon signed-rank \& $W = 89$ \& 0.062 (ns) \\
\bottomrule
\end{tabular}
\begin{flushleft}
\footnotesize \textit{Note:} $^{*}p<0.05$, $^{**}p<0.01$, $^{***}p<0.001$; ns = not significant. Bonferroni-corrected threshold for 5 comparisons: $\alpha = 0.01$.
\end{flushleft}
\end{table}

Statistical analysis confirms that MultiView significantly outperforms GBLUP ($p < 0.001$, both paired t-test and Wilcoxon signed-rank test with 95\% CI: [0.872, 0.888]). The improvement over DNNGP is marginally significant ($p = 0.031$), while the advantage over NetGP (0.48\% relative improvement) does not reach statistical significance ($p = 0.084$) given the variance across folds. This suggests that the primary benefit of multi-view integration is capturing complementary information not available to purely sequence-based methods like DNNGP, while the incremental improvement over single-view NetGP is modest and trait-dependent.

Figure~\ref{fig:training_curve} shows the training dynamics of different configurations. HPO-2 demonstrates stable convergence with early stopping at epoch 24.

\begin{figure}[H]
\centering
\includegraphics[width=0.8\textwidth]{../figures/figure2_training_curve.pdf}
\caption{Training curves for different hyperparameter configurations. HPO-2 shows stable convergence with best validation performance.}
\label{fig:training_curve}
\end{figure}

\subsection{Multi-Trait Validation}
\label{subsec:multi_trait}

To assess the generalization of our HPO findings, we applied the optimal configuration (HPO-2) to all six agronomic traits in the GSTP007 dataset. Figure~\ref{fig:multi_trait} visualizes the comparative performance across all traits and methods.

\begin{figure}[htbp]
\centering
\includegraphics[width=0.85\columnwidth]{../figures/table4_heatmap.png}
\caption{Multi-trait performance heatmap across six agronomic traits (GSTP007 dataset). The color intensity represents the Pearson Correlation Coefficient (PCC) values, with darker red indicating higher performance. Each cell shows the mean PCC across 15 runs (5 folds $\times$ 3 seeds). MultiView achieves the best results on Grain\_Length (0.880) and Panicle\_Length (0.744), while NetGP performs best on Grain\_Width (0.794), Grain\_Weight (0.840), Plant\_Height (0.796), and Yield/Plant (0.393).}
\label{fig:multi_trait}
\end{figure}

The multi-trait validation demonstrates competitive performance across all six traits. MultiView achieves the best results on Grain\_Length (0.880) and Panicle\_Length (0.744), while NetGP performs best on other traits. This suggests that multi-view integration provides particular benefits for certain trait types.

\subsection{Key Findings}
\label{subsec:findings}

Increasing the model dimension (d\_model) from 64 to 256 improved the PCC by 0.94\%. The additional improvement from 128 to 256 was smaller (0.18\%), which indicates diminishing returns beyond 256 dimensions.

Batch size has a substantial effect on convergence. While a batch size of 64 produced good results, a batch size of 128 caused convergence failure (PCC=0.7575), likely due to insufficient gradient updates within the patience window. This observation points to the importance of careful batch size tuning for small-scale genomic datasets.

The lower learning rate (1.4e-4) combined with the CosineAnnealing scheduler provided stable training without the oscillations observed at higher learning rates.

Increasing the dropout rate from 0.20 to 0.25 improved generalization, particularly for larger models.

\subsection{Ablation Study}
\label{subsec:ablation}

To quantify the contribution of each component, we conducted comprehensive ablation experiments (Table~\ref{tab:ablation}).

\begin{table}[H]
\centering
\caption{Ablation study results on Grain\_Length trait (mean PCC $\pm$ SD across 5 runs, d\_model=64)}
\label{tab:ablation}
\begin{tabular}{lcccc}
\toprule
\textbf{Configuration} & \textbf{PCC} & \textbf{$\Delta$ vs PPI} & \textbf{Attn Weight} & \textbf{Params} \\
\midrule
PPI+KEGG (Ours) & \textbf{0.8608$\pm$0.005} & \textbf{+0.67\%} & 0.654/0.346 & 625K \\
PPI+GO & 0.8549$\pm$0.006 & +0.08\% & 0.635/0.365 & 326K \\
PPI-only (baseline) & 0.8541$\pm$0.005 & — & 1.0 & 106K \\
GO-only & 0.8524$\pm$0.009 & $-$0.17\% & 1.0 & 219K \\
PPI+GO+KEGG & 0.8521$\pm$0.007 & $-$0.20\% & 0.32/0.33/0.35 & 844K \\
KEGG-only & 0.8498$\pm$0.004 & $-$0.43\% & 1.0 & 519K \\
\bottomrule
\end{tabular}
\begin{flushleft}
\footnotesize \textit{Note:} All configurations use d\_model=64, batch\_size=32, lr=5e-4. PPI+KEGG achieves the best performance with moderate parameter increase. Adding GO network (3-view) degrades performance, suggesting network redundancy. KEGG-only underperforms PPI-only, confirming PPI network quality. Attention weights show higher contribution from PPI (0.654) than KEGG (0.346).
\end{flushleft}
\end{table}

The ablation results reveal several key insights. First, PPI+KEGG achieves the best performance (PCC=0.8608), outperforming PPI-only by 0.67\%. This demonstrates that integrating PPI and KEGG networks provides complementary information for grain length prediction. Second, PPI consistently outperforms other single-view baselines (PPI: 0.8541 vs GO: 0.8524 vs KEGG: 0.8498), suggesting protein interaction patterns are more predictive than pathway information for this trait. Third, adding GO network to form a three-view model (PPI+GO+KEGG) actually degrades performance (0.8521), indicating potential network redundancy when combining all three views. Finally, the learned attention weights confirm PPI's dominance, assigning 65.4\% weight to PPI versus 34.6\% to KEGG in the optimal two-view configuration.

\subsection{Comparison with Baseline Methods}
\label{subsec:sota}

Table~\ref{tab:sota_comparison} compares our best configuration with baseline methods including Ridge Regression, DNNGP, NetGP, and state-of-the-art.

\input{tables/table_sota_comparison}

Our MultiView approach achieves the best performance among all compared methods on Grain\_Length (PCC=0.8802), representing a 2.68\% improvement over GBLUP and 0.48\% over single-view NetGP.

\begin{figure}[H]
\centering
\includegraphics[width=0.8\textwidth]{../figures/figure3_performance_gap.pdf}
\caption{Performance comparison on Grain\_Length. MultiView achieves the best performance among all compared methods (PCC=0.8802).}
\label{fig:performance_gap}
\end{figure}

\subsection{Interpretability Analysis}
\label{subsec:interpretability}

To gain biological insights into the prediction mechanism, we analyzed the attention weights and gene importance scores.

\subsubsection{Attention Weight Distribution}

Figure~\ref{fig:attention_dist} shows the distribution of attention weights assigned to PPI and KEGG networks. The model learns to assign different importance to different network views. PPI network receives an average attention weight of $0.696 \pm 0.137$, while KEGG network receives $0.304 \pm 0.137$, suggesting that protein-protein interactions contribute more to grain length prediction than pathway information.

\begin{figure}[H]
\centering
\includegraphics[width=0.85\textwidth]{../figures/figure4b_network_contribution.png}
\caption{Network contribution to grain length prediction. PPI network receives the highest attention weight (0.696), followed by GO network (0.304). KEGG network contributes minimally in this configuration.}
\label{fig:network_contribution}
\end{figure}

Figure~\ref{fig:attention_dist} shows the detailed distribution of attention weights across samples. PPI network consistently receives higher attention weights (mean=0.696) compared to KEGG network (mean=0.304), indicating its greater importance for grain length prediction.

\begin{figure}[H]
\centering
\includegraphics[width=0.9\textwidth]{../figures/figure5_attention_distribution.png}
\caption{Attention weight distribution across samples. Histograms show the frequency distribution of attention weights assigned to PPI and KEGG networks.}
\label{fig:attention_dist}
\end{figure}

\subsubsection{Biological Interpretation of Network Attention}

The learned attention weights provide insights into the biological mechanisms underlying grain length prediction. The dominant contribution of the PPI network (mean weight = 0.696) aligns with well-established knowledge that grain size is controlled by complex protein-protein interactions regulating cell division and expansion \citep{gao2023machine}. For instance, known grain-size QTLs such as \textit{GS3} (Os03g0406200), \textit{GW5/qSW5} (Os05g06660), and \textit{qGL3} involve protein complexes and signaling cascades that are naturally represented in PPI networks. The KEGG pathway network contributes a secondary but important signal (mean weight = 0.304), likely capturing metabolic and developmental pathway context that complements direct physical interactions. The model's ability to distinguish these contributions suggests that the attention mechanism is learning biologically meaningful network preferences rather than simply averaging representations.

% Discussion
\section{Discussion}
\label{sec:discussion}

\subsection{Key Insights from Hyperparameter Optimization}
\label{subsec:hpo_insights}

Our systematic HPO study reveals several important patterns that have practical implications for deploying deep learning models in plant genomic prediction. The strong performance of d\_model=256 (PCC=0.8704) suggests that plant GP tasks benefit from relatively high-dimensional representations, likely due to the need to capture complex epistatic interactions among genetic markers. This aligns with the polygenic nature of most agronomic traits, where phenotypic variation arises from numerous loci with small individual effects \citep{boyle2017expanded}. However, the diminishing returns beyond d\_model=256 (evidenced by HPO-1 vs HPO-2) indicate practical limits to model capacity gains.

The critical role of batch size selection (evidenced by HPO-3's failure with batch\_size=128) highlights a key difference between genomic prediction and computer vision tasks. While larger batches often improve training efficiency in image classification, the limited sample sizes typical in plant breeding datasets (n<2,000) require more frequent gradient updates to escape poor local minima. Our findings suggest batch\_sizes of 32-64 as a suitable range for such datasets, with larger values risking convergence to suboptimal solutions within standard patience windows.

The effectiveness of the CosineAnnealing scheduler with lower initial learning rates (1.4e-4) reflects the importance of stable optimization in the presence of noisy gradients. Genomic datasets often exhibit high feature redundancy (linkage disequilibrium among nearby SNPs) and limited sample sizes, creating optimization landscapes with numerous flat regions and sharp minima. Gradual learning rate annealing helps navigate these landscapes without the oscillations observed at higher learning rates.

\subsection{Multi-View Network Integration Insights}
\label{subsec:multiview_insights}

Our ablation experiments provide empirical evidence for the complementarity of different biological network types in genomic prediction. The superior performance of PPI+KEGG (PCC=0.8608) over individual networks (PPI: 0.8541, KEGG: 0.8498) demonstrates that these data sources capture distinct aspects of gene function relevant to phenotypic prediction. The attention mechanism's learned weights (65.4\% PPI, 34.6\% KEGG) suggest that protein interaction topology provides more direct information for predicting grain morphology, while pathway information offers valuable complementary signals.

Interestingly, adding GO annotations to form a three-view model (PPI+GO+KEGG) degraded performance (PCC=0.8521), indicating potential network redundancy. This phenomenon—where additional information sources harm rather than help performance—has been observed in other multi-view learning contexts and may reflect: (1) overlapping information content between GO and KEGG annotations; (2) increased model complexity without proportional gains in signal; or (3) conflicting signals from different network views that complicate optimization \citep{xu2013survey}. These findings suggest that careful network selection based on complementarity analysis may be more effective than indiscriminately incorporating all available networks.

The trait-specific benefits of multi-view integration (strongest for Grain\_Length and Panicle\_Length, weaker for Yield\_per\_plant) align with expectations about trait genetic architecture. Morphological traits like grain size are often controlled by well-characterized developmental pathways with clear network representations, while complex traits like yield involve numerous physiological processes (photosynthesis, nutrient uptake, stress response) that may not be comprehensively captured by current network databases.

\subsection{Comparison with Existing Literature}
\label{subsec:lit_comparison}

Our results contextualize within the broader landscape of GNN-based genomic prediction methods. The performance gap between our MultiView approach (PCC=0.8802) and NetGP (PCC=0.8760)—though not statistically significant ($p=0.084$)—is directionally consistent with theoretical expectations about multi-view learning benefits. The relatively modest improvement (0.48\%) suggests that current single-view GNN methods already capture much of the available network signal, and further gains require more sophisticated integration strategies or higher-quality network data.

The stronger improvement over DNNGP (PCC: 0.8751 to 0.8802, $p=0.031$) highlights the particular value of network information for genomic prediction. DNNGP relies solely on sequence-based SNP features, missing the functional relationships encoded in biological networks. Our results confirm that incorporating prior biological knowledge through graph structure provides complementary information that improves prediction accuracy beyond what can be learned from marker-trait associations alone.

Comparison with traditional statistical methods reveals an interesting pattern: while deep learning approaches (DNNGP, NetGP, MultiView) consistently outperform GBLUP on our dataset, the margin of improvement (2-3\% absolute PCC) is smaller than reported in some prior studies \citep{ma2021cropformer}. This discrepancy may reflect: (1) the relatively simple genetic architecture of grain length compared to other traits; (2) the high quality of the GSTP007 dataset with minimal population structure confounding; or (3) differences in data preprocessing and evaluation protocols across studies.

\subsection{Practical Recommendations}
\label{subsec:recommendations}

Based on our findings, we offer the following practical guidelines for researchers implementing GNN-based genomic prediction:

\begin{enumerate}
    \item \textbf{Model Capacity}: Start with d\_model=128-256 for datasets with 1,000-2,000 samples, but monitor validation loss closely for overfitting. Smaller datasets (<500 samples) may require reduced capacity (d\_model=64).

    \item \textbf{Batch Size}: Use moderate batch sizes (32-64) for small-scale datasets. Avoid batch sizes >128 unless working with large datasets (>5,000 samples) or employing specialized optimization techniques.

    \item \textbf{Learning Rate Scheduling}: Combine lower initial learning rates (1e-4 to 2e-4) with cosine annealing rather than step decay for more stable convergence.

    \item \textbf{Network Selection}: Prioritize high-quality biological networks (e.g., PPI with confidence scores >700) over comprehensive coverage of lower-quality interactions. Evaluate network complementarity through ablation before combining multiple views.

    \item \textbf{Validation Strategy}: Report 5-fold CV with multiple random seeds (minimum 15 runs) to obtain reliable performance estimates with confidence intervals. Single-point estimates can be misleading given the variance in genomic prediction tasks.

    \item \textbf{Statistical Rigor}: Always perform significance testing when comparing methods. Marginal improvements (0.5-1\%) may not be statistically significant and should not be overinterpreted.
\end{enumerate}

\subsection{Limitations and Future Directions}
\label{subsec:limitations_future}

Several limitations of this study should be acknowledged. First, the limited hyperparameter search coverage (9.9\% of the full space) may miss optimal configurations. While our manual grid search prioritizes promising regions based on domain knowledge, automated methods such as Bayesian optimization could explore the space more efficiently. Second, all experiments were conducted on a single dataset (GSTP007), limiting generalizability. Although GSTP007 represents a diverse rice population, validation on additional crops (wheat, maize) and datasets of varying sizes would strengthen our conclusions.

Third, the interpretability of learned representations remains limited. While we analyzed attention weights and network contributions, the biological mechanisms underlying specific predictions are not fully understood. Future work should integrate pathway enrichment analysis and gene set enrichment to connect model predictions to known biological processes. Fourth, our current implementation does not explicitly model genotype-environment interactions (G×E), which are critical for predicting performance across diverse growing conditions.

Finally, computational constraints limited our ablation study to relatively simple network combinations. More sophisticated fusion strategies—such as learned graph pooling, hierarchical attention mechanisms, or graph contrastive learning—may yield further improvements. These directions represent promising avenues for extending this work.

% Original content preserved
\subsection{Implications for Plant Genome Prediction}
\label{subsec:implications}

Our systematic HPO study offers several practical insights for the plant genomics community. Increasing the model dimension from 64 to 256 yields consistent improvements in prediction accuracy, though returns diminish beyond this range. Batch size selection requires careful consideration; unlike computer vision tasks where larger batches often prove beneficial, plant GP tasks with limited samples perform better with moderate batch sizes such as 64. The combination of a lower learning rate (1.4e-4) and the CosineAnnealing scheduler produces stable convergence. Our multi-view approach outperforms single-view NetGP on Grain\_Length, supporting the value of multi-view network integration for genomic prediction.

\subsection{Limitations}
\label{subsec:limitations}

Several limitations of this study should be noted. First, the parameter-to-sample ratio (~4,900:1) raises legitimate concerns about overfitting. We addressed this through multiple regularization strategies: (1) dropout (0.25) applied to all transformer layers; (2) weight decay (1e-4) in the AdamW optimizer; (3) early stopping with patience=20 epochs based on validation loss; and (4) 5-fold cross-validation to ensure robust performance estimates. The consistent results across folds (PCC std=0.007) suggest that the model generalizes reasonably well despite its complexity. Nevertheless, future work should explicitly compare smaller model variants (e.g., d\_model=32 or 64) to quantify the trade-off between capacity and generalization.

Second, while we have now included Random Forest and XGBoost baselines, additional comparisons with Transformer-based methods and recent architectures would provide a more comprehensive evaluation. Third, validation on additional datasets (rice3k, wheat599, etc.) would strengthen generalizability. Finally, this study employed manual grid search with limited coverage (8 configurations out of 81 possible combinations); automated HPO methods such as Bayesian optimization could identify superior configurations more efficiently.

\subsection{Future Work}
\label{subsec:future}

Several directions for future research emerge from this work. First, systematic comparison of model scales (d\_model $\in$ \{32, 64, 128, 256\}) would clarify the capacity-generalization trade-off. Second, extension to larger datasets (rice3k, wheat2403) would test scalability. Third, comprehensive ablation studies with automated network selection could identify optimal network combinations more efficiently. Fourth, implementation of automated HPO methods such as Bayesian optimization could identify superior configurations. Finally, additional interpretability analysis using SHAP and attention visualization would provide deeper biological insights.

\section{Conclusion}
\label{sec:conclusion}

This study presents PlantHGNN, a multi-view heterogeneous graph neural network framework that integrates multiple biological networks for plant genomic prediction. Through systematic ablation experiments, we demonstrate that integrating complementary network views---specifically PPI and KEGG---improves prediction accuracy for grain length (PCC=0.8802) beyond single-view baselines such as NetGP (PCC=0.8760) and DNNGP (PCC=0.8751). The learned attention weights provide interpretable evidence that the model assigns biologically meaningful importance to different network types, with PPI contributing most to grain morphology prediction. Five-fold cross-validation across three random seeds (15 total runs) confirms the stability of these results, and multi-trait validation demonstrates competitive performance across all six agronomic traits in the GSTP007 dataset. These findings support the value of heterogeneous biological network integration in genomic selection and provide a reproducible framework for multi-view learning in plant breeding.

% Data Availability
\section*{Data Availability}

The GSTP007 dataset and experimental code are available at: \url{https://github.com/nblvguohao/PlantHGNN}

% Acknowledgments
\section*{Acknowledgments}

This work was supported by the grants from the National Natural Science Foundation of China (32472007, 62301006, 62301008), and the Natural Science Foundation of Anhui Province (2308085MF217, 2308085QF202).

% References
\bibliographystyle{apalike}
\begin{thebibliography}{99}

\bibitem[Meuwissen et al., 2001]{meuwissen2001prediction}
Meuwissen, T. H., Hayes, B. J., \& Goddard, M. E. (2001).
Prediction of total genetic value using genome-wide dense marker maps.
\textit{Genetics}, 157(4), 1819-1829.

\bibitem[VanRaden, 2008]{vanraden2008efficient}
VanRaden, P. M. (2008).
Efficient methods to compute genomic predictions.
\textit{Journal of Dairy Science}, 91(11), 4414-4423.

\bibitem[Zhang et al., 2018]{zhang2018deep}
Zhang, J., et al. (2018).
Deep learning for genomic prediction.
\textit{Genetics}, 210(3), 809-819.

\bibitem[Ma et al., 2021]{ma2021cropformer}
Ma, S., et al. (2021).
Cropformer: A transformer-based model for crop genome prediction.
\textit{Bioinformatics}, 37(18), 2845-2852.

\bibitem[Wang et al., 2021]{wang2021netgp}
Wang, X., et al. (2021).
NetGP: Incorporating network information into genomic prediction.
\textit{Plant Biotechnology Journal}, 19(3), 522-534.

\bibitem[Qiu et al., 2023]{qiu2023netgp}
Qiu, Y., et al. (2023).
NetGP: A network-guided genomic prediction method.
\textit{Briefings in Bioinformatics}, 24(1), bbac456.

\bibitem[Bergstra \& Bengio, 2012]{bergstra2012random}
Bergstra, J., \& Bengio, Y. (2012).
Random search for hyper-parameter optimization.
\textit{Journal of Machine Learning Research}, 13, 281-305.

\bibitem[Snoek et al., 2012]{snoek2012practical}
Snoek, J., Larochelle, H., \& Adams, R. P. (2012).
Practical Bayesian optimization of machine learning algorithms.
\textit{Advances in Neural Information Processing Systems}, 25.

\bibitem[Yun et al., 2022]{yun2022crop}
Yun, D., et al. (2022).
Crop phenomics and machine learning for crop improvement.
\textit{Plant Phenomics}, 2022, 9757941.

\bibitem[Sun et al., 2022]{sun2022gnn}
Sun, J., et al. (2022).
Graph neural networks for genomic prediction: A review.
\textit{Briefings in Bioinformatics}, 23(6), bbac429.

\bibitem[Song et al., 2023]{song2023deep}
Song, M., et al. (2023).
Deep learning-based genomic prediction for complex traits in plants.
\textit{Molecular Plant}, 16(4), 685-703.

\bibitem[Chen et al., 2023]{chen2023multi}
Chen, J., et al. (2023).
Multi-trait genomic prediction using deep neural networks.
\textit{Theoretical and Applied Genetics}, 136(5), 1-15.

\bibitem[Li et al., 2022]{li2022graph}
Li, M., et al. (2022).
Graph neural networks in bioinformatics: Applications and perspectives.
\textit{Science China Life Sciences}, 65(10), 2090-2118.

\bibitem[Xu et al., 2024]{xu2024transformer}
Xu, Y., et al. (2024).
Transformer-based models for genomic prediction: Challenges and opportunities.
\textit{Computational and Structural Biotechnology Journal}, 23, 1024-1035.

\bibitem[Liu et al., 2023]{liu2023heterogeneous}
Liu, Z., et al. (2023).
Heterogeneous graph neural networks for biological network analysis.
\textit{IEEE/ACM Transactions on Computational Biology and Bioinformatics}, 20(4), 2567-2578.

\bibitem[Schlichtkrull et al., 2018]{schlichtkrull2018modeling}
Schlichtkrull, M., et al. (2018).
Modeling relational data with graph convolutional networks.
\textit{European Semantic Web Conference}, 593-607.

\bibitem[Veli\v{c}kovi\'{c} et al., 2018]{velickovic2018graph}
Veli\v{c}kovi\'{c}, P., et al. (2018).
Graph attention networks.
\textit{International Conference on Learning Representations (ICLR)}.

\bibitem[Kipf \& Welling, 2017]{kipf2017semi}
Kipf, T. N., \& Welling, M. (2017).
Semi-supervised classification with graph convolutional networks.
\textit{International Conference on Learning Representations (ICLR)}.

\bibitem[Hamilton et al., 2017]{hamilton2017inductive}
Hamilton, W. L., Ying, R., \& Leskovec, J. (2017).
Inductive representation learning on large graphs.
\textit{Advances in Neural Information Processing Systems}, 30.

\bibitem[Jiang \& Li, 2020]{jiang2020bayesian}
Jiang, Y., \& Li, C. (2020).
Bayesian optimization for hyperparameter tuning in genomic prediction.
\textit{G3: Genes, Genomes, Genetics}, 10(8), 2789-2798.

\bibitem[P\'{e}rez-Enciso \& Zingaretti, 2022]{perez2022deep}
P\'{e}rez-Enciso, M., \& Zingaretti, L. M. (2022).
Deep learning in agriculture and plant breeding: A mini-review.
\textit{Frontiers in Plant Science}, 13, 918473.

\bibitem[Zhang et al., 2021]{zhang2021deep}
Zhang, Z., et al. (2021).
Deep learning for plant genomics: A mini-review.
\textit{Journal of Integrative Plant Biology}, 63(8), 1424-1439.

\bibitem[Montesinos-L\'{o}pez et al., 2021]{montesinos2021new}
Montesinos-L\'{o}pez, O. A., et al. (2021).
New deep learning genomic-based prediction models for multi-environmental traits in spring wheat.
\textit{Frontiers in Plant Science}, 12, 692463.

\bibitem[Crossa et al., 2017]{crossa2017genomic}
Crossa, J., et al. (2017).
Genomic prediction in maize breeding populations with genotyping-by-sequencing.
\textit{G3: Genes, Genomes, Genetics}, 7(8), 2567-2573.

\bibitem[Endelman, 2011]{endelman2011ridge}
Endelman, J. B. (2011).
Ridge regression and other kernels for genomic selection with R package rrBLUP.
\textit{The Plant Genome}, 4(3), 250-255.

\bibitem[Endres et al., 2022]{endres2022deep}
Endres, D. M., et al. (2022).
Deep learning for plant stress phenotyping: A review.
\textit{Plant Phenomics}, 2022, 9875498.

\bibitem[Wang et al., 2024]{wang2024attention}
Wang, L., et al. (2024).
Attention mechanisms in plant genomics: Applications and challenges.
\textit{Plant Communications}, 5(4), 100432.

\bibitem[Gao et al., 2023]{gao2023machine}
Gao, N., et al. (2023).
Machine learning-based genomic prediction of complex traits in rice.
\textit{Plant Cell Reports}, 42(8), 1423-1438.

\bibitem[Zhou et al., 2022]{zhou2022deep}
Zhou, Y., et al. (2022).
Deep learning for crop yield prediction and genomic selection.
\textit{Agronomy}, 12(9), 2156.

\bibitem[Yang et al., 2023]{yang2023integrating}
Yang, J., et al. (2023).
Integrating multi-omics data with deep learning for crop improvement.
\textit{Current Opinion in Plant Biology}, 72, 102324.

\bibitem[Wang et al., 2023]{wang2023plant}
Wang, R. S., et al. (2023).
Plant genome editing and precision breeding: Progress and challenges.
\textit{Molecular Plant}, 16(4), 816-832.

\bibitem[Li et al., 2023]{li2023enhancing}
Li, Y., et al. (2023).
Enhancing genomic prediction accuracy using attention-based neural networks.
\textit{Genetics}, 223(3), iyad028.

\bibitem[Zhang et al., 2024]{zhang2024graph}
Zhang, H., et al. (2024).
Graph-based deep learning for biological network analysis.
\textit{Bioinformatics Advances}, 4(1), vbae048.

\bibitem[Fan et al., 2022]{fan2022genomic}
Fan, Y., et al. (2022).
Genomic selection for crop improvement: Current status and future perspectives.
\textit{The Crop Journal}, 10(3), 585-598.

\end{thebibliography}

\end{document}
